Optimal-pH prediction was formulated as a regression task. For an enzyme sequence \(x_i\), each model produced a scalar prediction \(\hat{y}_i\) for the experimentally measured optimal pH \(y_i\). Mean absolute error (MAE) was used as the primary metric because it is directly interpretable in pH units:

\begin{equation}
    \mathrm{MAE} = \frac{1}{n}\sum_{i=1}^{n} |y_i - \hat{y}_i|.
    \label{eq:mae}
\end{equation}

Additional metrics were root mean squared error (RMSE), coefficient of determination (\(R^2\)), Spearman correlation, and Pearson correlation. Bootstrap confidence intervals were estimated on test-set predictions using 500 bootstrap resamples.

Two feature representations were evaluated. The handcrafted representation included sequence length, amino-acid composition over the 20 standard residues, acidic and basic residue fractions, a simple net-charge proxy, hydrophobic, polar, aromatic, and aliphatic residue fractions, GRAVY hydrophobicity, and Biopython-derived physicochemical descriptors where available. The embedding representation used frozen ESM-2 embeddings from \texttt{facebook/esm2\_t6\_8M\_UR50D}. Sequence embeddings were generated by mean-pooling residue-token representations after excluding special and padding tokens. The resulting embedding matrix had shape 9855 by 320.

Generated evaluation strategies included a row-wise random split, an exact-duplicate-safe split, MMseqs2 sequence-cluster-controlled splits at 90\%, 70\%, and 50\% sequence identity, an organism-held-out split, and an EC-held-out split. The row-wise random split assigned records independently to train, validation, and test partitions. The grouped strategies assigned every sequence duplicate, sequence cluster, organism, or EC number to one partition only. MMseqs2 was run with the Windows release available for this analysis using `easy-cluster`, \texttt{--min-seq-id} 0.90, 0.70, or 0.50, \texttt{-c 0.80}, \texttt{--cov-mode 0}, and four threads; the corresponding commands are generated by \texttt{src/clustering.py}. These settings define cluster-disjoint partitions, but cluster membership alone does not prove that every test sequence is below the threshold against every training sequence. A nearest-training-sequence audit using \texttt{scripts/audit\_cluster\_similarity.py} is therefore required before interpreting these splits as homology exclusion. A dataset-provided split was also evaluated for continuity with the input file, but it is not treated as a leakage-controlled split.

The regressors were a mandatory mean predictor, a training-median predictor, ridge regression, random forest regression, and histogram gradient boosting regression. Hyperparameters were fixed before evaluation from \texttt{config/default.yaml}; no model or hyperparameter was selected using test-set performance. The validation partition was retained as a protocol check and for future tuning, but the reported winners are descriptive rankings by test MAE, not validation-selected final models. Training and validation records were not combined before fitting, so every reported test prediction comes from a model fitted on the designated training partition only. All model and split generation used seed 42; the random forest and histogram gradient boosting estimators received this seed, and bootstrap intervals used seed 42. Scalers and imputers were fitted only on the training partition inside scikit-learn pipelines.

Bootstrap intervals resampled individual test rows, so they quantify row-level uncertainty and may be optimistic when related proteins have dependent errors. Paired Wilcoxon comparisons likewise operate at the row level; Holm correction addresses multiplicity but not dependence. These limitations motivate reporting the paired MAE differences directly and treating the inferential results as exploratory.
